\section{Số trung bình và mốt của mẫu số liệu ghép nhóm}

\subsection{Đề 1}
\Opensolutionfile{ans}[ans/ans-11-C5B1-DE1-LC]
\caulc
\begin{ex}%%[1D5H1-1]
Dựa trên đồ thị thể hiện thi đánh giá năng lực của một trường đại học vào năm 2020 dưới đây. 
\begin{center}
	\begin{tikzpicture}[scale=1,>=stealth, font=\footnotesize, line join=round, line cap=round]
		\def\xmin{0} \def\xmax{13} % Hệ trục
		\def\ymin{0} \def\ymax{9} % Hệ trục
		\draw[->] (\xmin,0)--(\xmax,0)node[below]{(Điểm)};
		\draw[->] (0,\ymin)--(0,\ymax)node[above]{(Số học sinh)};
		\foreach \y/\z in {1/20, 2/40, 3/60, 4/80, 5/100, 6/120, 7/140}{
			\draw (0,\y)node[left]{$\z$};
		}
		\foreach \y/\z in {1/20, 2/40, 3/60, 4/80, 5/100, 6/120, 7/140}{
			\draw[gray!30](0,\y)--(13,\y);
		}
		\foreach \x in {0.25,0.75,1.25,1.75,2.25, 2.75, 3.25, 10.75, 11.25, 11.75}{
			\draw (\x,0)node[above]{$0$};}
		\foreach \x/\t in {0.25/[0;50], 0.75/[51;100] ,1.25/[101;150], 1.75/[151;200], 2.25/[201;250], 2.75/[251;300], 3.25/[301;350], 3.75/[351;400], 4.25/[401;450], 4.75/[451;500], 5.25/[501;550], 5.75/[551;600], 6.25/[601;650], 6.75/[651;700], 7.25/[701;750], 7.75/[751;800], 8.25/[801;850], 8.75/[851;900], 9.25/[901;950], 9.75/[951;1000], 10.25/[1001;1050], 10.75/[1051;1100], 11.25/[1101;1150], 11.75/[1151;1200]}{
			\draw (\x,-0.775)node[rotate=90]{$\t$};}
		\draw[fill=gray] (3.5,0) rectangle (4,0.05) (3.75,0.05)node[above]{$1$};
		\draw[fill=gray] (4,0) rectangle (4.5,0.05*8) (4.25,0.05*8)node[above]{$8$};
		\draw[fill=gray] (4.5,0) rectangle (5,0.05*24) (4.75,0.05*24)node[above]{$24$};
		\draw[fill=gray] (5,0) rectangle (5.5,0.05*54) (5.25,0.05*54)node[above]{$54$};
		\draw[fill=gray] (5.5,0) rectangle (6,0.05*95) (5.75,0.05*95)node[above]{$95$};
		\draw[fill=gray] (6,0) rectangle (6.5,0.05*95) (6.25,0.05*95)node[above]{$95$};
		\draw[fill=gray] (6.5,0) rectangle (7,0.05*133) (6.75,0.05*133)node[above]{$133$};
		\draw[fill=gray] (7,0) rectangle (7.5,0.05*122) (7.25,0.05*122)node[above]{$122$};
		\draw[fill=gray] (7.5,0) rectangle (8,0.05*104) (7.75,0.05*104)node[above]{$104$};
		\draw[fill=gray] (8,0) rectangle (8.5,0.05*62) (8.25,0.05*62)node[above]{$62$};
		\draw[fill=gray] (8.5,0) rectangle (9,0.05*55) (8.75,0.05*55)node[above]{$55$};
		\draw[fill=gray] (9,0) rectangle (9.5,0.05*21) (9.25,0.05*21)node[above]{$21$};
		\draw[fill=gray] (9.5,0) rectangle (10,0.05*12) (9.75,0.05*12)node[above]{$12$};
		\draw[fill=gray] (10,0) rectangle (10.5,0.05*1) (10.25,0.05*1)node[above]{$1$};
		\draw (\xmax/2,\ymax)node[above]{\textbf{Điểm thi đánh giá năng lực của một trường đại học năm 2020}};
	\end{tikzpicture}
\end{center}
Tổng số học sinh tham gia kì thi đánh giá năng lực trên là
	\choice
	{$780$}
	{\True $787$}
	{$696$}
	{$697$}
	\loigiai{
		Tổng số học sinh là $1+8+24+54+95+95+133+122+104+62+55+21+12+1=787$ (học sinh).
	}
\end{ex}

\begin{ex}%%[1D5H1-1]
Dựa trên đồ thị thể hiện thi đánh giá năng lực của một trường đại học vào năm 2020 dưới đây. 
\begin{center}
	\begin{tikzpicture}[scale=1,>=stealth, font=\footnotesize, line join=round, line cap=round]
		\def\xmin{0} \def\xmax{13} % Hệ trục
		\def\ymin{0} \def\ymax{9} % Hệ trục
		\draw[->] (\xmin,0)--(\xmax,0)node[below]{(Điểm)};
		\draw[->] (0,\ymin)--(0,\ymax)node[above]{(Số học sinh)};
		\foreach \y/\z in {1/20, 2/40, 3/60, 4/80, 5/100, 6/120, 7/140}{
			\draw (0,\y)node[left]{$\z$};
		}
		\foreach \y/\z in {1/20, 2/40, 3/60, 4/80, 5/100, 6/120, 7/140}{
			\draw[gray!30](0,\y)--(13,\y);
		}
		\foreach \x in {0.25,0.75,1.25,1.75,2.25, 2.75, 3.25, 10.75, 11.25, 11.75}{
			\draw (\x,0)node[above]{$0$};}
		\foreach \x/\t in {0.25/[0;50], 0.75/[51;100] ,1.25/[101;150], 1.75/[151;200], 2.25/[201;250], 2.75/[251;300], 3.25/[301;350], 3.75/[351;400], 4.25/[401;450], 4.75/[451;500], 5.25/[501;550], 5.75/[551;600], 6.25/[601;650], 6.75/[651;700], 7.25/[701;750], 7.75/[751;800], 8.25/[801;850], 8.75/[851;900], 9.25/[901;950], 9.75/[951;1000], 10.25/[1001;1050], 10.75/[1051;1100], 11.25/[1101;1150], 11.75/[1151;1200]}{
			\draw (\x,-0.775)node[rotate=90]{$\t$};}
		\draw[fill=gray] (3.5,0) rectangle (4,0.05) (3.75,0.05)node[above]{$1$};
		\draw[fill=gray] (4,0) rectangle (4.5,0.05*8) (4.25,0.05*8)node[above]{$8$};
		\draw[fill=gray] (4.5,0) rectangle (5,0.05*24) (4.75,0.05*24)node[above]{$24$};
		\draw[fill=gray] (5,0) rectangle (5.5,0.05*54) (5.25,0.05*54)node[above]{$54$};
		\draw[fill=gray] (5.5,0) rectangle (6,0.05*95) (5.75,0.05*95)node[above]{$95$};
		\draw[fill=gray] (6,0) rectangle (6.5,0.05*95) (6.25,0.05*95)node[above]{$95$};
		\draw[fill=gray] (6.5,0) rectangle (7,0.05*133) (6.75,0.05*133)node[above]{$133$};
		\draw[fill=gray] (7,0) rectangle (7.5,0.05*122) (7.25,0.05*122)node[above]{$122$};
		\draw[fill=gray] (7.5,0) rectangle (8,0.05*104) (7.75,0.05*104)node[above]{$104$};
		\draw[fill=gray] (8,0) rectangle (8.5,0.05*62) (8.25,0.05*62)node[above]{$62$};
		\draw[fill=gray] (8.5,0) rectangle (9,0.05*55) (8.75,0.05*55)node[above]{$55$};
		\draw[fill=gray] (9,0) rectangle (9.5,0.05*21) (9.25,0.05*21)node[above]{$21$};
		\draw[fill=gray] (9.5,0) rectangle (10,0.05*12) (9.75,0.05*12)node[above]{$12$};
		\draw[fill=gray] (10,0) rectangle (10.5,0.05*1) (10.25,0.05*1)node[above]{$1$};
		\draw (\xmax/2,\ymax)node[above]{\textbf{Điểm thi đánh giá năng lực của một trường đại học năm 2020}};
	\end{tikzpicture}
\end{center}
	Giá trị đại diện cho nhóm chứa mốt của mẫu số liệu ghép nhóm trên là
	\choice
	{$625{,}5$}
	{\True $675{,}5$}
	{$725{,}5$}
	{$775{,}5$}
	\loigiai{
		Nhóm có số học sinh nhiều nhất là $\left[651;700\right)$.
	}
\end{ex}

\begin{ex}%[1D5N1-1]
	Dựa trên bảng số liệu về chiều cao của $100$ học sinh một trường trung học phổ thông dưới đây.
	\begin{center}
	\begin{tabular}{|c|c|c|}
		\hline
		Nhóm & Chiều cao (cm) & số học sinh \\
		\hline
		$1$ & $[150;153)$ & $7$ \\
		\hline
		$2$ & $[153;156)$ & $13$ \\
		\hline
		$3$ & $[156;159)$ & $40$ \\
		\hline
		$4$ & $[159;162)$ & $21$ \\
		\hline
		$5$ & $[162;165)$ & $13$ \\
		\hline
		$6$ & $[165;168)$ & $6$ \\
		\hline
	\end{tabular}
\end{center}
	$160{,}5$ là giá trị đại diện cho nhóm
	\choice
	{$2$}
	{$3$}
	{\True $4$}
	{$5$}
	\loigiai{
		$160{,}5$ là giá trị đại diện cho nhóm $4$ vì $160{,}5=\dfrac{159+162}{2}$.
	}
\end{ex}

\begin{ex}%[1D5H1-3]
	Số trung bình của mẫu số liệu trên thuộc khoảng nào trong các khoảng dưới đây?
	\choice
	{$[7;9)$}
	{\True $[9;11)$}
	{$[11;13)$}
	{$[13;15)$}
	\loigiai{
		Bảng tần số ghép nhóm theo giá trị đại diện là
		\begin{center}
			\begin{tabular}{|c|c|c|c|c|c|}
				\hline Doanh thu &{$[5; 7)$} &{$[7; 9)$} &{$[9; 11)$} &{$[11; 13)$} &{$[13; 15)$} \\
				\hline Giá trị đại diện &$6$ &$8$ &$10$ &$12$ &$14$ \\
				\hline Số ngày & 2 & 7 & 7 & 3 & 1 \\
				\hline
			\end{tabular}
		\end{center}
		Số trung bình: $\overline{x}=\dfrac{2\cdot 6+7\cdot 8+7\cdot 10+3\cdot 12+1\cdot 14}{20}=9{,}4$.
	}
\end{ex}

\begin{ex}%[1D5H1-4]
	Mốt của mẫu số liệu trên thuộc khoảng nào trong các khoảng dưới đây?
	\choice
	{$[7;9)$}
	{\True $[9;11)$}
	{$[11;13)$}
	{$[13;15)$}
	\loigiai{
		Có 2 nhóm chứa mốt của mẫu số liệu trên, do đó
		\begin{itemize}
			\item Xét nhóm $[7;9)$:
			\[M_0=7+\dfrac{7-2}{(7-2)+(7-7)}\cdot (9-7)=9. \]
			\item Xét nhóm $[9;11)$:
			\[M_0^\prime=7+\dfrac{7-7}{(7-7)+(7-3)}\cdot (11-9)=9. \]
		\end{itemize}
		Do đó Mốt của mẫu số liệu bằng $9$.
	}
\end{ex}

\begin{ex}%[1D5H1-1]
Khảo sát thời gian tập thể dục trong ngày của một số học sinh khối 11 thu được mẫu số liệu ghép nhóm sau
\begin{center}
	\begin{tabular}{|c|c|c|c|c|c|}
		\hline 
		Thời gian (phút)&$[0;20)$&$[20;40)$&$[40;60)$&$[60;80)$&$[80;100)$\\
		\hline 
		Số học sinh &$5$&$9$&$12$&$10$&$6$\\
		\hline 
	\end{tabular}
\end{center}
Giá trị đại diện của nhóm $[20; 40)$ là
\choice 
{$10$}
{$20$}
{\True $30$}
{$40$}
\loigiai{
	Trong mỗi khoảng thời gian, giá trị đại diện là trung bình cộng của giá trị hai đầu mút nên ta có giá trị đại diện của nhóm $[20;40)$ là $\dfrac{20+40}{2}=30$.
} 
\end{ex}
\begin{ex}%[1D5H1-4]
Khảo sát thời gian tập thể dục trong ngày của một số học sinh khối 11 thu được mẫu số liệu ghép nhóm sau
\begin{center}
	\begin{tabular}{|c|c|c|c|c|c|}
		\hline 
		Thời gian (phút)&$[0;20)$&$[20;40)$&$[40;60)$&$[60;80)$&$[80;100)$\\
		\hline 
		Số học sinh &$5$&$9$&$12$&$10$&$6$\\
		\hline 
	\end{tabular}
\end{center}
Nhóm chứa mốt của mẫu số liệu này là 
\choice
{$[20;40)$}
{\True $[40;60)$}
{$[60;80)$}
{$[80;100)$}
\loigiai{Tần số lớn nhất là $12$ nên nhóm chứa mốt là nhóm $[40;60)$.} 
\end{ex}

\begin{ex}%[1D5H1-4]
	Người ta tiến hành phỏng vấn $40$ người về một mẫu áo sơ mi mới. Người điều tra yêu cầu cho điểm mẫu áo đó theo thang điểm là $100$. Kết quả được trình bày trong \textit{Bảng dưới}.
	\begin{center}
	\begin{tabular}{|c|c|c|}
		\hline
		Nhóm & Tần số & Tần số tích lũy \\
		\hline
		[50;60) & 4 & 4 \\
		\hline
		[60;70) & 5 & 9 \\
		\hline
		[70;80) & 23 & 32 \\
		\hline
		[80;90) & 6 & 38 \\
		\hline
		[90;100) & 2 & 40 \\
		\hline
		\multicolumn{3}{|c|}{$n=40$} \\
		\hline
	\end{tabular}
	\end{center}
	Mốt của mẫu số liệu ghép nhóm trên (làm tròn đến kết quả hàng đơn vị) là
	\choice
	{$73$}
	{$74$}
	{\True $75$}
	{$76$}
	\loigiai{
	Ta thấy nhóm $3$ ứng với nửa khoảng $[70;80)$ là nhóm có tần số lớn nhất với $u=70;g=10;n_3=23$. Nhóm $2$ có tần số $n_2=5$, nhóm $4$ có tần số $n_4=6$.\\
	Áp dụng công thức, ta có mốt của mẫu số liệu là:
	$$ M_o=70+\left(\dfrac{23-5}{2\cdot23-5-6}\right)\cdot10 \approx 75.$$}
\end{ex}

\begin{ex}%[1D5H1-4]
Bảng số liệu ghép nhóm sau cho biết chiều cao (cm) của $50$ học sinh lớp $11A$.
\begin{center}
	\begin{tabular}{|c|c|c|c|c|c|c|}
		\hline
		Khoảng chiều cao (cm)	& $\left[145;150 \right)$ & $\left[150;155 \right)$ & $\left[155;160 \right)$ & $\left[160;165 \right)$&$\left[165;170 \right)$  \\
		\hline
		Số học sinh&$7$	& $14$ & $10$ &$10$  & $9$ \\
		\hline
	\end{tabular}
\end{center}
Tính mốt của mẫu số liệu ghép nhóm này (làm tròn đến hàng đơn vị).
	\choice
	{$151$}
	{$152$}
	{\True $153$}
	{$154$}
	\loigiai{
Tần số lớn nhất là $14$ nên nhóm chứa mốt là nhóm $\left[150;155 \right)$. Ta có $j=2$, $a_2=150$, $m_2=14$, $m_1=7$, $m_3=10$, $h=5$. Do đó $$M_o=150+\dfrac{14-7}{\left(14-7\right)+\left(14-10\right)\cdot 5}\approx 153.$$	
	}
\end{ex}

\begin{ex}%[1D5H1-4]
	Một công ty xây dựng khảo sát khách hàng xem họ có nhu cầu mua nhà ở mức giá nào. Kết quả khảo sát được ghi lại ở bảng sau
	\begin{center}
		\begin{tabular}{|c|c|c|c|c|c|}
			\hline \begin{tabular}{c}
				\textbf{Mức giá} \\
				\textbf{(triệu đồng/$\mathrm{m}^2)$}
			\end{tabular} &{$[10; 14)$} &{$[14; 18)$} &{$[18; 22)$} &{$[22; 26)$} &{$[26; 30)$} \\
			\hline \textbf{Số khách hàng} & 54 & 78 & 120 & 45 & 12 \\
			\hline
		\end{tabular}
	\end{center}
	Tìm mốt của mẫu số liệu ghép nhóm trên (làm tròn đến hàng phần mười).
	\choice
	{$18{,}3$}
	{\True $19{,}4$}
	{$12{,}4$}
	{$15{,}5$}
	\loigiai{
		Nhóm chứa mốt của mẫu số liệu trên là nhóm $[18; 22)$.\\ Do đó $u_m=18$, $n_{m-1}=78$, $n_m=120$, $n_{m+1}=45$, $u_{m+1}-u_m=22-18=4$.\\
		Mốt của mẫu số liệu ghép nhóm là
		\[M_0=18+\dfrac{120-78}{(120-78)+(120-45)} \cdot 4=\dfrac{758}{39} \approx 19{,}4. \]
	}
\end{ex}

\begin{ex}%[1D5H1-3]
Diện tích các tỉnh và thành phố khu vực Nam Bộ được thống kê ở bảng sau
\begin{center}
\begin{tabular}{|c|c|c|c|}
\hline Tỉnh/ thành phố & Diện tích $\left(\mathbf{k m}^2\right)$ & Tỉnh/ thành phố & Diện tích $\left.\mathbf{( k m}^{\mathbf{2}}\right)$ \\
\hline Bình Phước & 6877 & Vĩnh Long & 1526 \\
\hline Tây Ninh & 4041 & Đồng Tháp & 3384 \\
\hline Bình Dương & 2695 & An Giang & 3537 \\
\hline Đồng Nai & 5864 & Kiên Giang & 6349 \\
\hline Bà Rịa - Vũng Tàu & 1981 & Cần Thơ & 1439 \\
\hline TP.Hồ Chí Minh & 2061 & Hậu Giang & 1622 \\
\hline Long An & 4495 & Sóc Trăng & 3312 \\
\hline Tiền Giang & 2511 & Bạc Liêu & 2669 \\
\hline Bến Tre & 2395 & Cà Mau & 5221 \\
\hline Trà Vinh & 2358 & & \\
\hline
\end{tabular}
\newline
(Nguồn: Tổng cục Thống kê)
\end{center}
	Hãy ước lượng số trung bình của mẫu số liệu ghép nhóm trên (làm tròn đến hàng phần trăm).
	\choice
		{\True $3407{,}89$}
		{$3321{,}35$}
		{$3514{,}26$}
		{$3612{,}15$}
	\loigiai{
		Bảng tần số ghép nhóm bao gồm giá trị đại diện của các nhóm như sau
		\begin{center}
		\begin{tabular}{|l|c|c|c|c|}
		\hline Diện tích $\left(\mathrm{km}^2\right)$ & {$[1000 ; 2500)$} & {$[2500 ; 4000)$} & {$[4000 ; 5500)$} & {$[5500 ; 7000)$} \\
		\hline Giá trị đại diện & 1750 & 3250 & 4750 & 6250 \\
		\hline Số tỉnh/ thành phố & 7 & 6 & 3 & 3 \\
		\hline
		\end{tabular}
		\end{center}
		Khi đó, số trung bình của mẫu số liệu ghép nhóm là
		$$\bar{x}=\dfrac{7\cdot 1750+6 \cdot 3250+3 \cdot 4750+3 \cdot 6250}{19} \approx 3407{,}89.$$
}
\end{ex}

\begin{ex}%[1D5N1-1]
	Thời gian hoàn thành bài kiểm tra Toán $45$ phút của các bạn trong lớp được cho như sau
	\begin{center}
	\begin{tabular}{|c|c|c|c|c|c|}
	\hline
	Thời gian (phút) & $[25 ; 30)$ & {$[30 ; 35)$} & {$[35 ; 40)$} & {$[40 ; 45]$} \\
	\hline
	Số học sinh & $2$ & $7$ & $10$ & $25$ \\
	\hline
	\end{tabular}
	\end{center}
	Có bao nhiêu học sinh hoàn thành bài kiểm tra trước khi hết giờ trên 5 phút?
	\choice
	{\True $19$}
	{$18$}
	{$20$}
	{$17$}
	\loigiai{
	\begin{enumerate}
	\item Các nhóm số liệu $[25; 30)$, $[30; 35)$, $[35; 40)$, $[40; 45]$ với tần số tương ứng là $2,7,10,25$.
	\item Số học sinh hoàn thành bài kiểm tra trước khi hết giờ ít nhất $5$ phút là $2+7+10=19$.
	\end{enumerate}
	}
\end{ex}
\Closesolutionfile{ans}
\indapan{6}{ans/ans-11-C5B1-DE1-LC}

\cauds
\Opensolutionfile{ans}[ans/ans-11-C5B1-DE1-DS]
\begin{ex}%[1D1H5-3]
	Kết quả khảo sát cân nặng của $25$ quả cam ở mỗi lô hàng ${A, B}$ được cho ở bảng sau
	\begin{center}
	\begin{tabular}{|c|c|c|c|c|c|}
	\hline
	Cân nặng (gam)           & $[150 ; 155)$ & $[155 ; 160)$ & $[160 ; 165)$ & $[165 ; 170)$ & $[170 ; 175)$ \\ \hline
	Số quả cam ở lô hàng $A$ & 2               & 6               & 12              & 4               & 1               \\ \hline
	Số quả cam ở lô hàng $B$ & 1               & 3               & 7               & 10              & 4               \\ \hline
	\end{tabular}
	\end{center}
	\choiceTF
	{\True Giá trị đại diện nhóm $[150 ; 155)$ bằng $152{,}5$}
	{Cân nặng trung bình của mỗi quả cam ở lô $A$ là $163{,}7$ (gam)}
	{Cân nặng trung bình của mỗi quả cam ở lô $B$ là $162{,}1$ (gam)}
	{\True Theo số trung bình thì cam ở lô hàng ${B}$ nặng hơn cam ở lô hàng $A$}
	\loigiai{
	\begin{itemchoice}
	\itemch \textbf{Đúng}. Bảng thống kê số lượng cam theo giá trị đại diện
	\begin{center}
	\begin{tabular}{|c|c|c|c|c|c|}
	\hline
	Cân nặng (gam)           & $152{,}5$ & $157{,}5$ & $162{,}5$ & $167{,}5$ & $172{,}5$ \\ \hline
	Số quả cam ở lô hàng $A$ & 2               & 6               & 12              & 4               & 1               \\ \hline
	Số quả cam ở lô hàng $B$ & 1               & 3               & 7               & 10              & 4               \\ \hline
	\end{tabular}
	\end{center}	
	\itemch \textbf{Sai}. Cân nặng trung bình của mỗi quả cam ở lô $A$ là
	$${\bar{x}_A=\dfrac{152{,}5 \cdot 2+157{,}5 \cdot 6+162{,}5 \cdot 12+167{,}5 \cdot 4+172{,}5 \cdot 1}{25}=161{,}7 \text { (gam).}
	}$$
	\itemch \textbf{Sai}. Cân nặng trung bình của mỗi quả cam ở lô $B$ là
	$${\bar{x}_B=\dfrac{152,5 \cdot 1+157,5 \cdot 3+162,5 \cdot 7+167,5 \cdot 10+172,5 \cdot 4}{25}=165,1 \text { (gam).}
	}$$
	\itemch \textbf{Đúng}.  Ta thấy ${\bar{x}_A<\bar{x}_B}$. Vậy nếu so sánh theo số trung bình thì cam ở lô hàng $B$ nặng hơn cam ở lô hàng $A$.
	\end{itemchoice}
}
\end{ex}

\begin{ex}%[1D1H5-3]
	Số lượng người đi xem một bộ phim mới theo độ tuổi trong một rạp chiếu phim (sau $1$ h đầu công chiếu) được ghi lại theo bảng phân phối ghép nhóm sau:
	\begin{center}
	\begin{tabular}{|c|c|c|c|c|c|}
	\hline
	Độ tuổi  & $[10 ; 20)$ & $[20 ; 30)$ & $[30 ; 40)$ & $[40 ; 50)$ & $[50 ; 60)$ \\ \hline
	Số người & 6           & 12          & 16          & 7           & 2           \\ \hline
	\end{tabular}
	\end{center}
	\choiceTF
	{\True Giá trị đại diện nhóm $[50 ; 60)$ là $55$} 
	{\True Độ tuổi được dự báo là ít xem phim đó nhất là thuộc nhóm $[50;60)$}
	{\True Nhóm chứa mốt là nửa khoảng $[30 ; 40)$}
	{Độ tuổi được dự báo là thích xem phim đó nhiều nhất là 31 tuổi}
	\loigiai{
		\begin{itemchoice}
		\itemch \textbf{Đúng}. Giá trị đại diện nhóm $[50 ; 60)$ là $55$.
		\itemch \textbf{Đúng}. Độ tuổi được dự báo là ít xem phim đó nhất là thuộc nhóm $[50;60)$.
		\itemch \textbf{Đúng}. Nhóm chứa mốt là nửa khoảng $[30 ; 40)$.
		\itemch \textbf{Sai}. Ta có mốt là $M_0=30+\dfrac{16-12}{(16-12)+(16-7)} \cdot 10=\dfrac{430}{13} \approx 33{,}08$. Vậy độ tuổi được dự báo là thích xem phim đó nhiều nhất là $33$ tuổi.
		\end{itemchoice}
		}
\end{ex}

\begin{ex}%[1D1H5-3]
	Một nhà thực vật học đo chiều dài trung bình của $74$ lá cây (đơn vị: milimét) và thu được bảng tần số ghép nhóm như sau
	{\fontsize{10pt}{12pt}\selectfont
	\begin{center}	
	\begin{tabular}{|c|c|c|c|c|c|c|c|}
	\hline
	Nhóm             & $[5{,}45 ; 5{,}85)$ & $[5{,}85 ; 6{,}25)$ & $[6{,}25 ; 6{,}65)$ & $[6{,}65 ; 7{,}05)$ & $[7{,}05 ; 7{,}45)$ & $[7{,}45 ; 7{,}85)$ & $[7{,}85 ; 8{,}25)$ \\ \hline
	Giá trị đại diện & 5,65            & 6,05            & 6,45            & 6,85            & 7,25            & 7,65            & 8,05            \\ \hline
	Tần số           & 5               & 9               & 15              & 15              & 19              & 8               & 2               \\ \hline
	\end{tabular}
	\end{center}}
	\choiceTF
	{Chiều dài trung bình của $74$ lá cây bằng $\approx 6{,}4$ (mm)}
	{\True Độ dài nhóm là $0{,}4$}
	{Nhóm chứa mốt là $[7{,}05;7{,}45)$}
	{Mốt của mẫu số liệu ghép nhóm là $\approx 6{,}65$}
	\loigiai{
		\begin{itemchoice}
		\itemch \textbf{Sai}. Chiều dài trung bình của 74 lá cây là:
		$$\begin{aligned}
		\bar{x} & =\dfrac{5,65.5+6,05.9+6,45 \cdot 15+6,85 \cdot 19+7,25.16+7,65 \cdot 8+8,05.2}{74} \\
		& =\dfrac{5029}{740} \approx 6{,}8({~mm}) .
		\end{aligned}
		$$
		\itemch \textbf{Đúng}. Độ dài nhóm là $0{,}4$.
		\itemch \textbf{Sai}.Nhóm chứa mốt là ${[6,65 ; 7,05)}$,\\
		trong đó
		${u_m=6{,}65 ; u_{m+1}=7{,}05 ; u_{m+1}-u_m=0{,}4 ; n_m=19 ; n_{m-1}=15 ; n_{m+1}=16
		}$.
		\itemch \textbf{Sai}. Mốt của mẫu số liệu ghép nhóm là:
		${M_0=6{,}65+\dfrac{19-15}{(19-15)+(19-16)} \cdot 0{,}4 \approx 6{,}88.}$
		\end{itemchoice}
}
\end{ex}

\begin{ex}%[1D1H5-3]
	Số cuộc điện thoại một người thực hiện mỗi ngày trong 30 ngày được lựa chọn ngẫu nhiên được thống kê trong bảng sau:
	\begin{center}
	\begin{tabular}{|c|c|c|c|c|c|}
	\hline
	Số cuộc gọi  & $[2{,}5 ; 5{,}5)$ & $[5{,}5 ; 8{,}5)$ & $[8{,}5 ; 11{,}5)$ & $[11{,}5 ; 14{,}5)$ & $[14{,}5 ; 17{,}5)$ \\ \hline
	Số ngày & 5           & 13          & 7          & 3           & 2           \\ \hline
	\end{tabular}
	\end{center}
	\choiceTF
	{Số cuộc gọi trung bình mỗi ngày là $8{,}1$} 
	{\True Nhóm chứa mốt là $[5{,}5 ; 8{,}5)$}
	{\True Mốt của mẫu số liệu ghép nhóm là $\approx 7{,}21$}
	{Người đó thực hiện tối đa khoảng 8 cuộc gọi mỗi ngày}
	\loigiai{
		\begin{itemchoice}
		\itemch \textbf{Sai}. Ta viết lại bảng tần số ghép nhóm theo giá trị đại diện là
		\begin{center}
		\begin{tabular}{|c|c|c|c|c|c|}
		\hline
		Số cuộc gọi  & $[2{,}5 ; 5{,}5)$ & $[5{,}5 ; 8{,}5)$ & $[8{,}5 ; 11{,}5)$ & $[11{,}5 ; 14{,}5)$ & $[14{,}5 ; 17{,}5)$ \\ \hline
		Giá trị đại diện & 4           & 7          & 10          & 13           & 16 \\ \hline
		Số ngày & 5           & 13          & 7          & 3           & 2 \\ \hline
		\end{tabular}
		\end{center}
		Số cuộc gọi trung bình mỗi ngày là:
		$$\bar{x}=\dfrac{4\cdot5+7\cdot13+10\cdot7+13\cdot3+16\cdot2}{30}=8{,}4.$$
		\itemch \textbf{Đúng}. Nhóm chứa mốt là $[5{,}5; 8{,}5)$.
		\itemch \textbf{Đúng}. Ta có ${u_m=5,5 ; u_{m+1}=8,5 \Rightarrow u_{m+1}-u_m=3 ; n_m=13 ; n_{m-1}=5 ; n_{m+1}=7}$.\\
		Vì vậy mốt của mẫu số liệu ghép nhóm là
		$$M_0=5,5+\frac{13-5}{(13-5)+(13-7)} \cdot 3=\frac{101}{14} \approx 7{,}21.$$
		\itemch \textbf{Sai}. Vậy người đó thực hiện tối đa khoảng 7 cuộc gọi mỗi ngày.
		\end{itemchoice}
		}
\end{ex}
\Closesolutionfile{ans}
\indapan{2}{ans/ans-11-C5B1-DE1-DS}

\caukq
\Opensolutionfile{ans}[ans/ans-11-C5B1-DE1-KQ]
\begin{ex}%[1D5H1-4]
Tuổi thọ (năm) của 50 bình ắc quy ô tô được cho như sau
\begin{center}
	\begin{tabular}{|c|c|c|c|c|c|c|}
		\hline
		Tuổi thọ (năm)	& $\left[2;2{,}5 \right)$ & $\left[2{,}5;3 \right)$ & $\left[3;3{,}5 \right)$&$\left[3{,}5;4 \right)$&$\left[4;4{,}5 \right)$&$\left[4{,}5;5 \right)$  \\
		\hline
	Tần số &$4$	& $9$ & $14$ &$11$  & $7$&$5$ \\
		\hline
	\end{tabular}
\end{center}
	Xác định mốt (làm tròn đến hàng phần mười).
	\shortans{$3{,}3$}
	\loigiai{
	Tần số lớn nhất là $14$ nên nhóm chứa mốt là nhóm $\left[3;3, 5\right)$.\\
	Ta có $j=3$, $a_3=3$, $m_3=14$, $m_2=9$, $m_4=11$, $h=0$, $5$. Do đó
	$$M_o=3+\dfrac{14-9}{(14-9)+(14-11)}\cdot 0{,}5\approx 3{,}3.$$	
}
\end{ex}

\begin{ex}%[1D5H1-3]
	Thống kê điểm trung bình môn Toán của một số học sinh lớp $11$ được cho ở bảng sau
	\begin{center}
		\begin{tabular}{|c|c|c|c|c|c|c|c|}
			\hline Khoảng điểm &{$[6{,}5; 7)$} &{$[7; 7{,}5)$} &{$[7{,}5; 8)$} &{$[8; 8{,}5)$} &{$[8{,}5;9)$} &{$[9; 9{,}5)$} &{$[9{,}5; 10)$} \\
			\hline Tần số & $8$ & $10$ & $16$ & $24$ & $13$ & $7$ & $4$ \\
			\hline
		\end{tabular}
	\end{center}
	Hãy ước lượng số trung bình của mẫu số liệu ghép nhóm trên (làm tròn đến hàng phần trăm).
	\shortans{$8{,}12$}
	\loigiai{
		Bảng tần số ghép nhóm theo giá trị đại diện
		\begin{center}
			\begin{tabular}{|c|c|c|c|c|c|c|c|}
				\hline Khoảng điểm &{$[6{,}5; 7)$} &{$[7; 7{,}5)$} &{$[7{,}5; 8)$} &{$[8; 8{,}5)$} &{$[8{,}5;9)$} &{$[9; 9{,}5)$} &{$[9{,}5; 10)$} \\
				\hline Giá trị đại diện &$6{,}75$ &$7{,}25$ &$7{,}75$ &$8{,}25$ &$8{,}75$ &$9{,}25$ &$9{,}75$ \\
				\hline Tần số & $8$ & $10$ & $16$ & $24$ & $13$ & $7$ & $4$ \\
				\hline
			\end{tabular}
		\end{center}
		Điểm trung bình môn Toán của một số học sinh lớp $11$ là
		\[\overline{x}=\dfrac{8\cdot 6{,}75+10\cdot 7{,}25+16\cdot 7{,}75+24\cdot 8{,}25+13\cdot 8{,}75+7\cdot 9{,}25+4\cdot 9{,}75}{82}\approx 8{,}12.\]
	}
\end{ex}

\begin{ex}%[1D5H1-3]
	Để kiểm tra thời gian sử dụng pin của một chiếc điện thoại mới, chị An thống kê thời gian sử dụng điện thoại của mình từ lúc sạc đầy pin cho đến khi hết pin ở bảng sau:
	\begin{center}
		\begin{tabular}{|c|c|c|c|c|c|} \hline
			\text{Thời gian sử dụng (giờ)} & $\left[7;9\right)$ & $\left[9;11\right)$ & $\left[11;13\right)$ & $\left[13,15\right)$ & $\left[15;17\right)$ \\ \hline
			\text{Số lần} & 2 & 5 & 7 & 6 & 3 \\ \hline
		\end{tabular}
	\end{center}
	Hãy ước lượng thời gian sử dụng trung bình từ lúc chị An sạc đầy pin điện thoại cho tới khi hết pin (làm tròn đến hàng phần mười).
	\shortans{$12{,}3$}
	\loigiai{
		Thời gian sử dụng trung bình
			$$\overline{x}=\dfrac{2\cdot8 +5\cdot10 +7\cdot12 +6\cdot14 +3\cdot16}{2+5+7+6+3}=\dfrac{282}{23}\approx 12{,}3.$$
	}
\end{ex}

\begin{ex}%[1D5H1-4]
	Bảng sau thống kê số ca nhiễm mới SARS-CoV-2 mỗi ngày trong tháng 12/2021 tại Việt Nam.
	\begin{center}
		\begin{tabular}{|c|c|c|c|c|c|c|c|} \hline
			{Ngày} & {Số ca} & {Ngày} & {Số ca} & {Ngày} & {Số ca} & {Ngày} & \text{Số ca}  \\ \hline
			1 & 15\,139 &  9 & 15\,965 & 17 & 15\,685 & 25 & 16\,046 \\\hline
			2 & 14\,295 & 10 & 15\,474 & 18 & 16\,363 & 26 & 15\,667 \\\hline
			3 & 14\,254 & 11 & 16\,830 & 19 & 16\,586 & 27 & 15\,310 \\\hline
			4 & 14\,598 & 12 & 15\,264 & 20 & 15\,420 & 28 & 14\,866 \\\hline
			5 & 14\,927 & 13 & 16\,035 & 21 & 16\,806 & 29 & 14\,299 \\\hline
			6 & 15\,215 & 14 & 15\,871 & 22 & 17\,044 & 30 & 20\,454 \\\hline
			7 & 14\,433 & 15 & 16\,192 & 23 & 16\,860 & 31 & 17\,004 \\\hline
			8 & 15\,223 & 16 & 15\,720 & 24 & 16\,633 &&\\ \hline
		\end{tabular}
		\begin{flushright}
			(\textit{Nguồn}: worldmeters.info)
		\end{flushright}
	\end{center}
	Tìm mốt của mẫu số liệu ghép nhóm trên.
	\shortans{$15{,}7$}
	\loigiai{
		\textbf{\color{magenta}Bảng tần số ghép nhóm}
			\begin{center}
				\begin{tabular}{|c|c|c|c|c|c|} \hline
					Số ca (nghìn) & $\left[14;15{,}5\right)$ & $\left[15{,}5;17\right)$ & $\left[17;18{,}5\right)$ & $\left[18{,}5;20\right)$ & $\left[20;21{,}5\right)$ \\ \hline
					Số ngày & 13 & 15 & 2 & 0 & 1\\ \hline
				\end{tabular}
			\end{center}
			Mốt $M_o$ chứa trong nhóm $\left[u_{m};u_{m+1}\right)$
			$$M_o=u_{m}+\dfrac{n_{m}-n_{m-1}}{\left(n_{m}-n_{m-1}\right)+\left(n_{m}-n_{m+1}\right)}\left(u_{m+1}-u_{m}\right)$$
			Nhóm chứa mốt của mẫu số liệu trên là nhóm $\left[120;175\right)$\\
			Do đó $\begin{aligned}[t]
				& u_{m}=15{,}5 ; u_{m+1}=17  \Rightarrow u_{m+1}-u_{m}=17-15{,}5=1{,}5.\\
				& n_{m-1}=13 ; n_{m}=15 ; n_{m+1}=2. \\
				& M_o=15{,}5+\dfrac{15-13}{(15-13)+(15-2)}\cdot(17-15{,}5) =\dfrac{157}{10} = 15{,}7.
			\end{aligned}$
	}
\end{ex}

\begin{ex}%[1D5H1-4]
	Số cuộc gọi điện thoại một nguời thực hiện mỗi ngày trong $30$ ngày được lựa chọn ngẫu nhiên được thống kê trong bảng sau:
	\begin{center}
		\begin{tabular}{|c|c|c|c|c|c|}
			\hline Số cuộc gọi &{$[3; 5]$} &{$[6; 8]$} &{$[9; 11]$} &{$[12; 14]$} &{$[15; 17]$} \\
			\hline Số ngày & 5 & 13 & 7 & 3 & 2 \\
			\hline
		\end{tabular}
	\end{center}
	Tìm mốt của mẫu số liệu ghép nhóm trên (làm tròn đến hàng phần mười).
	\shortans{$7{,}2$}
	\loigiai{
		Do số cuộc gọi là số nguyên nên ta hiệu chỉnh lại như sau:
		\begin{center}
			\begin{tabular}{|c|c|c|c|c|c|}
				\hline Số cuộc gọi &{$[2{,}5; 5{,}5)$} &{$[5{,}5; 8{,}5)$} &{$[8{,}5; 11{,}5)$} &{$[11{,}5; 14{,}5)$} &{$[14{,}5; 17{,}5)$} \\
				\hline Số ngày & 5 & 13 & 7 & 3 & 2 \\
				\hline
			\end{tabular}
		\end{center}
		Nhóm chứa mốt của mẫu số liệu trên là nhóm $[5{,}5; 8{,}5)$.\\
			Do đó $u_m=5{,}5$; $n_{m-1}=5$; $n_m=13$; $n_{m+1}=7$; $u_{m+1}-u_m=8{,}5-5{,}5=3$.\\
			Mốt của mẫu số liệu ghép nhóm là
			\[M_0=5{,}5+\dfrac{13-5}{(13-5)+(13-7)} \cdot 3=\dfrac{101}{14} \approx 7{,}2. \]
	}
\end{ex}

\begin{ex}%[1D5H1-3]
	Người ta đếm số xe ô tô đi qua một trạm thu phí mỗi phút trong khoảng thời gian từ $9$ giờ đến $9$ giờ $30$ phút sáng. Kết quả được ghi lại ở bảng sau
	\begin{center}
		\begin{tabular}{|c|c|c|c|c|c|c|c|c|c|c|c|c|c|c|}
			\hline $15$ & $16$ & $13$ & $21$ & $17$ & $23$ & $15$ & $21$ & $6$ & $11$ & $12$ & $23$ & $19$ & $25$ & $11$ \\
			\hline $25$ & $7$ & $29$ & $10$ & $28$ & $29$ & $24$ & $6$ & $11$ & $23$ & $11$ & $21$ & $9$ & $27$ & $15$ \\
			\hline
		\end{tabular}
	\end{center}
	Hãy ước lượng trung bình số xe đi qua trạm thu phí trong mỗi phút (làm tròn đến hàng phần mười).
	\shortans{$17{,}7$}
	\loigiai{
		Bảng tần số ghép nhóm
			\begin{center}
				\begin{tabular}{|c|c|c|c|c|c|}
					\hline Số xe &{$[6; 10]$} &{$[11; 15]$} &{$[16; 20]$} &{$[21; 25]$} &{$[26; 30]$} \\
					\hline Số lần & $5$ & $9$ & $3$ & $9$ & $4$ \\
					\hline
				\end{tabular}
			\end{center}
		Bảng tần số ghép nhóm (theo giá trị đại diện) được hiệu chỉnh lại như sau
			\begin{center}
				\begin{tabular}{|c|c|c|c|c|c|}
					\hline Số xe &{$[5{,}5; 10{,}5)$} &{$[10{,}5; 15{,}5)$} &{$[15{,}5; 20{,}5)$} &{$[20{,}5; 25{,}5)$} &{$[25{,}5; 30{,}5)$} \\
					\hline Giá trị đại diện &{$8$} &{$13$} &{$18$} &{$23$} &{$28$} \\
					\hline Số lần & $5$ & $9$ & $3$ & $9$ & $4$ \\
					\hline
				\end{tabular}
			\end{center}
			Số xe trung bình đi qua trạm qua bảng tần số ghép nhóm là
			\[\overline{x}=\dfrac{8\cdot 5+13\cdot 9+18\cdot 3+23\cdot 9+28\cdot 4}{30}\approx 17{,}7\ (\text{xe}). \]
	}
\end{ex}
\Closesolutionfile{ans}
\indapan{6}{ans/ans-11-C5B1-DE1-KQ}
\subsection{Đề 2}
\Opensolutionfile{ans}[ans/ans-11-C5B1-DE2-LC]
\caulc
\begin{ex}%[1D5H1-1]
	Nhóm số liệu rời rạc $k_1-k_2$ với $k_1$, $k_2\in \mathbb{N}$, $k_1<k_2$ là nhóm gồm các giá trị
	\choice
	{$k_1$ và $k_2$}
	{$k+1$, $\ldots$, $k_2$}
	{$k_1$, $\ldots$, $k_2+1$}
	{\True $k_1$, $k_1+1$, $\ldots$, $k_2$}
\loigiai{
Với $k_1$, $k_2\in \mathbb{N}$ và $k_1<k_2$ suy ra nhóm số liệu rời rạc từ $k_1$ đến $k_2$ gồm $k_1$, $k_1+1$, $k_1+2$, $\ldots$, $k_2$.
}
\end{ex}

\begin{ex}%[1D5H1-1]
	Giá trị đại diện của nhóm $[a_i;a_{i+1})$ là
	\choice
	{$a_i$}
	{$a_{i+1}$}
	{$\dfrac{a_{i+1}-a_i}{2}$}
	{\True $\dfrac{a_{i+1}+a_i}{2}$}
\loigiai{
Giá trị đại diện của nhóm $[a_i;a_{i+1})$ là $\dfrac{a_{i+1}+a_i}{2}$.
}
\end{ex}

\begin{ex}%[1D5H1-4]
	Mẫu số liệu ghép nhóm với tần số các nhóm bằng nhau có số mốt là
	\choice
	{\True $0$}
	{$1$}
	{$2$}
	{$3$}
\loigiai{
Khi tần số của các nhóm số liệu bằng nhau thì mẫu số liệu ghép nhóm không có mốt.
}
\end{ex}

\begin{ex}%[1D5H1-3]
	Cho mẫy số liệu ghép nhóm về tuổi thọ (đơn vị tính là năm) của một loại bóng đèn mới như sau.
	\begin{center}
		\begin{longtable}{|c|c|c|c|c|}
			\hline
			Tuổi thọ (năm) & $[2;3{,}5)$ & $[3{,}5;5)$ & $[5;6{,}5)$ & $[6{,}5;8)$\\
			\hline
			Số bóng đèn & $8$ & $22$ & $35$ & $15$\\
			\hline
		\end{longtable}
	\end{center}
	Số trung bình của mẫu số liệu là
	\choice
	{$5{,}0$}
	{\True $5{,}32$}
	{$5{,}75$}
	{$6{,}5$}
\loigiai{
Ta có bảng sau
\begin{center}
	\begin{longtable}{|c|c|c|c|c|}
		\hline
		Tuổi thọ (năm) & $[2;3{,}5)$ & $[3{,}5;5)$ & $[5;6{,}5)$ & $[6{,}5;8)$\\
		\hline
		Giá trị đại diện & $2{,}75$ & $4{,}25$ & $5{,}75$ & $7{,}25$\\
		Số bóng đèn & $8$ & $22$ & $35$ & $15$\\
		\hline
	\end{longtable}
\end{center}
Tổng số bóng đèn là $n=80$.\\
Số trung bình của mẫu số liệu là
$$\overline{x}=\dfrac{2{,}75\cdot 8+4{,}25\cdot 22+5{,}75\cdot  35+7{,}25\cdot 15}{80}=5{,}31875\approx 5{,}32.$$
}
\end{ex}

\begin{ex}%[1D5H1-2]
	Cho mẫy số liệu ghép nhóm về tuổi thọ (đơn vị tính là năm) của một loại bóng đèn mới như sau.
	\begin{center}
		\begin{longtable}{|c|c|c|c|c|}
			\hline
			Tuổi thọ (năm) & $[2;3{,}5)$ & $[3{,}5;5)$ & $[5;6{,}5)$ & $[6{,}5;8)$\\
			\hline
			Số bóng đèn & $8$ & $22$ & $35$ & $15$\\
			\hline
		\end{longtable}
	\end{center}
	Nhóm chứa mốt của mẫu số liệu là
	\choice
	{$[2;3{,}5)$}
	{$[3{,}5;5)$}
	{\True $[5;6{,}5)$}
	{$[6{,}5;8)$}
\loigiai{
Tần số lớn nhất là $35$ nên nhóm chứa mốt là nhóm $[5;6{,}5)$. 
}
\end{ex}

\begin{ex}%[1D5H1-2]
	Cho mẫy số liệu ghép nhóm về tuổi thọ (đơn vị tính là năm) của một loại bóng đèn mới như sau.
	\begin{center}
		\begin{longtable}{|c|c|c|c|c|}
			\hline
			Tuổi thọ (năm) & $[2;3{,}5)$ & $[3{,}5;5)$ & $[5;6{,}5)$ & $[6{,}5;8)$\\
			\hline
			Số bóng đèn & $8$ & $22$ & $35$ & $15$\\
			\hline
		\end{longtable}
	\end{center}
	Số mốt của mẫu số liệu ghép nhóm này là
	\choice
	{$0$}
	{\True $1$}
	{$2$}
	{$3$}
\loigiai{
Tần số lớn nhất là $35$ nên nhóm chứa mốt là nhóm $[3{,}5;5)$. Ta có $a_3=3{,}5$, $m_3=35$, $m_2=22$, $m_4=15$, $h=1{,}5$.\\
Do đó mốt của mẫu số liệu ghép nhóm này là
$$M_0=3{,}5+\dfrac{22-35}{(22-35)+(22-15)}\cdot 1{,}5=6{,}75.$$
Vậy mẫu số liệu ghép nhóm này có $1$ mốt.
}
\end{ex}

\begin{ex}%[1D5H1-2]
Khảo sát thời gian tập thể dục trong ngày của một số học sinh khối 11 thu được mẫu số liệu ghép nhóm sau
\begin{center}
	\begin{tabular}{|c|c|c|c|c|c|}
		\hline 
		Thời gian (phút)&$[0;20)$&$[20;40)$&$[40;60)$&$[60;80)$&$[80;100)$\\
		\hline 
		Số học sinh &$5$&$9$&$12$&$10$&$6$\\
		\hline 
	\end{tabular}
\end{center}
	Mẫu số liệu ghép nhóm này có số mốt là
	\choice 
	{$50$}
	{\True $52$}
	{$29$}
	{$35$}
	\loigiai{ 
	Vì mẫu số liệu chỉ có một lớp có tần số lớn nhất, nên mẫu này có $1$ mốt.\\
	Tần số lớn nhất là $12$ nên nhóm chứa mốt là nhóm thứ ba $[40; 60)$.\\
	Ta có $j=3$; $a_3=40$; $m_3=12$; $m_2=9$; $m_4=10$; $ h=20$.\\
	 Khi đó $M_O=a_j+\dfrac{m_j-m_{j-1}}{\left(m_j-m_{j-1}\right)+\left(m_j-m_{j+1}\right)}\cdot h=40+\dfrac{12-9}{(12-9)+(12-10)}\cdot 20=52$.
}
\end{ex}

\begin{ex}%[1D5H1-2]
	Thống kê chỉ số chất lượng không khí $(AQI)$ tại một địa điểm vào các ngày trong tháng $6 / 2022$ được cho trong bảng sau:
	\begin{center}
	\begin{tabular}{|c|c|c|c|c|c|}
	\hline Chỉ số $AQI$ & {$[0; 50)$} & {$[50; 100)$} & {$[100; 150)$} & {$[150; 200)$} & Trên 200 \\
	\hline Số ngày & $5$ & $11$ & $7$ & $4$ & $3$ \\
	\hline
	\end{tabular}
	\end{center}
	Chất lượng không khí được xem là tốt nếu $AQI$ nhỏ hơn $50$, là trung bình nếu $AQI$ từ $50$ đến dưới $100$. Trong tháng $6/2022$ tại địa điểm này có bao nhiêu ngày chất lượng không khí dưới mức trung bình?
	\choice 
	{$15$}
	{\True $14$}
	{$17$}
	{$12$}
	\loigiai{
		\begin{enumerate}
		\item Trong tháng $6/2022$ có $5$ ngày chỉ số $AQI$ dưới $50$; $11$ ngày chỉ số $AQI$ từ $50$ đến dưới $100$; $7$ ngày chỉ số $AQI$ từ $100$ đến dưới $150$; $4$ ngày chỉ số $AQI$ từ $150$ đến dưới $200$; $3$ ngày chỉ số $AQI$ trên $200$.
		\item Số ngày chất lượng không khí dưới mức trung bình là $7+4+3=14$.
		\end{enumerate}
	}
\end{ex}

\begin{ex}%[1D5H1-3]
	Quãng đường (km) các cầu thủ (không tính thủ môn) chạy trong một trận bóng đá tại giải ngoại hạng Anh được cho trong bảng thống kê sau:
	\begin{center}
		\begin{tabular}{|c|c|c|c|c|c|}
			\hline Quãng đường & {$[2 ; 4)$} & {$[4 ; 6)$} & {$[6 ; 8)$} & {$[8 ; 10)$} & {$[10 ; 12)$} \\
			\hline Số cầu thủ & 2 & 5 & 6 & 9 & 3 \\
			\hline
		\end{tabular}
	\end{center}
		Tính quãng đường trung bình một cầu thủ chạy trong trận đấu này.
		\choice 
		{$7{,}25$}
		{$8{,}23$}
		{\True $7{,}48$}
		{$7{,}67$}
		\loigiai{
			Tổng số cầu thủ là $n=2+5+6+9+3=25$.\\
			Quãng đường trung bình một cầu thủ chạy trong trận đấu này là
			$$\overline{x}=\dfrac{2\cdot 3+5\cdot 5+6\cdot 7+9\cdot 9+3\cdot 11}{25}=7{,}48~(\mathrm{km}).$$ 
		}
\end{ex}
	
\begin{ex}%[1D5H1-3]
	Số nguyện vọng đăng kí vào đại học của các bạn trong lớp được thống kê trong bảng sau
	\begin{center}
		\begin{longtable}{|c|c|c|c|c|}
			\hline
			Số nguyện vọng & $1-3$ & $4-6$ & $7-9$ & $10-12$\\
			\hline
			Số học sinh & $5$ & $18$ & $13$ & $7$\\
			\hline
		\end{longtable}
	\end{center}
	Trung bình một bạn trong lớp đăng kí bao nhiêu nguyện vọng.
	\choice 
	{$7{,}45$}
	{$6{,}67$}
	{$7{,}42$}
	{\True $6{,}53$}	
	\loigiai{
	Số trung bình của mẫu số liệu là
	$$\overline{x}=\dfrac{5\cdot 2+18\cdot 5+13\cdot 8+7\cdot 11}{43}\approx 6{,}53.$$
}
\end{ex}

\begin{ex}%[1D5H1-3]
	Nồng độ cồn trong hơi thở (đơn vị tính là miligam/1 lít khí thở) của $20$ lái xe ô tô vi phạm được cho như sau
	\begin{center}
		\begin{longtable}{|c|c|c|c|c|c|c|c|c|c|}
			\hline
			$0{,}09$ & $0{,}18$ & $0{,}47$ & $1{,}20$ & $0{,}28$ & $0{,}45$ & $0{,}72$ & $0{,}15$ & $0{,}75$ & $0{,}36$\\
			\hline
			$0{,}21$ & $0{,}15$ & $0{,}23$ & $0{,}30$ & $0{,}41$ & $0{,}13$ & $0{,}05$ & $0{,}38$ & $0{,}42$ & $0{,}79$\\
			\hline
		\end{longtable}
	\end{center}
Theo quy định, mức phạt nồng độ cồn đối với lái xe ô tô như sau:
\begin{description}
	\item[Mức 1.] Nồng độ cồn trong hơi thở chưa vượt quá $0{,}25$ phạt từ $6$ đến $8$ triệu đồng;
	\item[Mức 2.] Nồng độ cồn trong hơi thở từ trên $0{,}25$ đến $0{,}4$ phạt từ $16$ đến $18$ triệu đồng;
	\item[Mức 3.] Nồng độ cồn trong hơi thở vượt quá $0{,}4$ phạt từ $30$ đến $40$ triệu đồng.
\end{description}
	Trung bình mỗi lái xe bị phạt bao nhiêu tiền?
	\choice 
	{$23{,}5$}
	{$22{,}1$}
	{$24{,}4$}
	{\True $20{,}2$}		
	\loigiai{
	Ta có bảng số liệu ghép nhóm như sau
	\begin{center}
		\begin{longtable}{|c|c|c|c|}
			\hline
			Tiền phạt (triệu đồng) & $6-8$ & $16-18$ & $30-40$\\
			\hline
			Người vi phạm & $8$ & $4$ & $8$\\
			\hline 
		\end{longtable}
	\end{center}
	Cỡ mẫu $n=20$. Số tiền trung bình một người bị phạt là 
	$$\overline{x}=\dfrac{8\cdot 7+4\cdot 17+8\cdot  35}{20}=20{,}2\text{ (triệu đồng).}$$
}
\end{ex}

\begin{ex}%[1D5H1-3]
	$100$ người thực hiện bài trắc nghiệm để đo chỉ số $IQ$, kết quả thu được như sau:
	\begin{center}
		\begin{tabular}{|c|c|c|c|c|c|c|}
			\hline
			Chỉ số $IQ$	& Dưới $70$ & $\left[70;85\right)$ & $\left[85;115\right)$ & $\left[115;130\right)$ & $\left[130;145\right)$ & Từ $145$ trở lên \\
			\hline
			Số người & $2$ & $15$ & $45$ & $20$ & $15$ & $3$ \\
			\hline
		\end{tabular}
	\end{center}
	Người có chỉ số $IQ$ từ $85$ đến dưới $115$ là ở mức trung bình. Xác định tỉ lệ người có $IQ$ cao hơn mức trung bình.
	\choice 
	{$41\%$}
	{$35\%$}
	{$40\%$}
	{\True $38\%$}		
	\loigiai{
	\begin{listEX}
	\item Các nhóm số liệu gồm Dưới $70$; $\left[70;85\right); \left[85;115\right); \left[115;130\right); \left[130;145\right)$;\\
	 Từ $145$ trở lên với các tần số tương ứng là $2, 15, 45, 20, 15, 3$.
	\item Số người có chỉ số $IQ$ cao hơn mức trung bình là $20+15+3=38$.\\
	Vậy tỉ lệ người có chỉ số $IQ$ cao hơn mức trung bình là $\dfrac{38}{100}=38\%$. 
	\end{listEX}}
\end{ex}

\Closesolutionfile{ans}
\indapan{6}{ans/ans-11-C5B1-DE2-LC}
\cauds
\Opensolutionfile{ans}[ans/ans-11-C5B1-DE2-DS]
\begin{ex}%[1D1H5-3]
	Số câu trả lời đúng một bài thi trắc nghiệm môn Sinh học gồm 50 câu của lớp 11A ở một trường THPT như sau
	\begin{center}
	\begin{tabular}{|c|c|c|c|c|c|}
	\hline
	Thâm niên (Số năm)  & $[14;21)$ & $[21; 28)$ & $[28;35)$ & $[35; 42)$ & $[42;49)$ \\ \hline
	Số giáo viên & 4           & 8          & 25          & 6           & 7           \\ \hline
	\end{tabular}
	\end{center}
	\choiceTF
	{\True Giá trị đại diện của nhóm $[14; 21)$ là $17{,}5$} 
	{\True Giá trị đại diện của nhóm $[21; 28)$ là $24{,}5$}
	{\True Giá trị đại diện của nhóm $[42; 49)$ là $45{,}5$}
	{Số câu đúng trung bình là $32{,}26$}
	\loigiai{
		\begin{itemchoice}
		\itemch \textbf{Đúng}. Giá trị đại diện của nhóm $[14; 21)$ là $17{,}5$.
		\itemch \textbf{Đúng}.Giá trị đại diện của nhóm $[21; 28)$ là $24{,}5$.
		\itemch \textbf{Đúng}. Giá trị đại diện của nhóm $[42; 49)$ là $45{,}5$.
		\itemch \textbf{Sai}. 
		\begin{center}
		\begin{tabular}{|c|c|c|c|c|c|}
		\hline
		Thâm niên (Số năm)  & $[14;21)$ & $[21; 28)$ & $[28;35)$ & $[35; 42)$ & $[42;49)$ \\ \hline
		Giá trị đại diện & $17{,}5$           & $24{,}5$          & $31{,}5$         & $38{,}5$           & $45{,}5$          \\ \hline
		Số giáo viên & 4           & 8          & 25          & 6           & 7           \\ \hline
		\end{tabular}
		\end{center}
		Số câu đúng trung bình
		$$\bar{x}=\dfrac{4\cdot17{,}5+8\cdot24{,}5+25\cdot31{,}5+6\cdot38{,}5+7\cdot45{,}5}{50}=32{,}06.$$
		\end{itemchoice}
		}
\end{ex}

\begin{ex}%[1D1H5-3]
	Thâm niên giảng dạy của một số giáo viên trường THPT được ghi lại ở bảng sau
	\begin{center}
	\begin{tabular}{|c|c|c|c|c|c|}
	\hline
	Thâm niên (Số năm)  & $[1;5)$ & $[5; 10)$ & $[10;15)$ & $[15; 20)$ & $[20;25)$ \\ \hline
	Số giáo viên & 4           & 12          & 16          & 8           & 3           \\ \hline
	\end{tabular}
	\end{center}
	\choiceTF
	{\True Cỡ mẫu của mẫu số liệu bằng $50$} 
	{\True Số trung bình của mẫu ghép nhóm là $11{,}84$}
	{\True Nhóm chứa mốt của mẫu số liệu trên là nhóm $[10 ; 15)$}
	{Mốt của mẫu số liệu ghép nhóm bằng $11{,}74$}
	\loigiai{
		\begin{itemchoice}
		\itemch \textbf{Đúng}. Ta viết lại bảng tần số ghép nhóm theo giá trị đại diện là
		\begin{center}
		\begin{tabular}{|c|c|c|c|c|c|}
		\hline
		Thâm niên (Số năm)  & $[1;5)$ & $[5; 10)$ & $[10;15)$ & $[15; 20)$ & $[20;25)$ \\ \hline
		Giá trị đại diện & 3           & $7{,}5$          & $12{,}5$          & $17{,}5$           & $22{,}5$           \\ \hline
		Số giáo viên & 4           & 12          & 16          & 8           & 3           \\ \hline
		\end{tabular}
		\end{center}
		Cỡ mẫu của mẫu số liệu bằng $50$.
		\itemch \textbf{Đúng}. Ước lượng số trung bình của mẫu ghép là
		$$\bar{x}=\dfrac{3\cdot4+7{,}5\cdot12+12{,}5\cdot16+17{,}5\cdot8+22{,}5\cdot3}{50}=11{,}84.$$
		\itemch \textbf{Đúng}. Nhóm chứa mốt của mẫu số liệu trên là nhóm $[10; 15)$.
		\itemch \textbf{Sai}. Do đó, $u_m=10, n_{m-1}=12, n_m=16, n_{m+1}=8, n_{m+1}-u_m=15-10=5$.
		$$M_{\circ}=10+\frac{16-12}{(16-12)+(16-8)} \cdot 5=11{,}66.$$
		\end{itemchoice}
		}
\end{ex}

\begin{ex}%[1D1H5-3]
	Một nhà nghiên cứu ghi lại thời gian (giờ) sử dụng Facbook của 30 học sinh trong 02 tuần.
	Kết quả thu được mẫu số liệu như sau:
	\begin{center}
	\begin{tabular}{cccccccccc}
	21 & 17 & 22 & 18 & 20 & 17 & 15 & 13 & 15 & 20 \\
	15 & 12 & 18 & 17 & 25 & 17 & 21 & 15 & 12 & 18 \\
	16 & 23 & 14 & 18 & 19 & 13 & 16 & 19 & 18 & 17
	\end{tabular}
	\end{center}
	\choiceTF
	{Số giờ trung bình của học sinh trong 02 tuần là $16{,}37$ giờ}
	{\True Tổng hợp kết quả thời gian sử dụng Facbook của học sinh vào bảng tần số ghép nhóm theo mẫu sau\\
	\begin{tabular}{|c|c|c|c|c|c|}
			\hline
			Số giờ  & $[12;15)$ & $[15; 18)$ & $[18;21)$ & $[21; 24)$ & $[24;27)$ \\ \hline
			Giá trị đại diện & $13{,}5$           & $16{,}5$          & $19{,}5$          & $22{,}5$           & $25{,}5$           \\ \hline
			Số học sinh & 5           & 12          & 8          & 4           & 1           \\ \hline
			\end{tabular}}
	{\True Nhóm chứa mốt của mẫu số liệu ý b) là nhóm $[15 ; 18)$}
	{\True Mốt của mẫu số liệu ý b) bằng $16{,}91$}
	\loigiai{
		\begin{itemchoice}
		\itemch \textbf{Sai}. Tổng số thời gian sử dụng Facbook của 30 học sinh là 521 giờ.\\
		Số giờ trung bình của học sinh trong 02 tuần là $\bar{x}=\frac{521}{30}=17{,}37$ giờ.
		\itemch \textbf{Đúng}. 
		\begin{center}
		\begin{tabular}{|c|c|c|c|c|c|}
					\hline
					Số giờ  & $[12;15)$ & $[15; 18)$ & $[18;21)$ & $[21; 24)$ & $[24;27)$ \\ \hline
					Giá trị đại diện & $13{,}5$           & $16{,}5$          & $19{,}5$          & $22{,}5$           & $25{,}5$           \\ \hline
					Số học sinh & 5           & 12          & 8          & 4           & 1           \\ \hline
			\end{tabular}
		\end{center}
		\itemch \textbf{Đúng}. Nhóm chứa mốt của mẫu số liệu trên là nhóm $[15 ; 18)$.
		\itemch \textbf{Đúng}. Do đó, $u_m=15 ; n_{m-1}=5 ; n_m=12 ; n_{m+1}=8 ; u_{m+1}-u_m=18-15=3,0$.\\
		$M_{\circ}=15+\frac{12-5}{(12-5)+(12-8)} \cdot 3=16{,}91.$
		\end{itemchoice}
}
\end{ex}

\begin{ex}%[1D1H5-3]
	Người ta tiến hành phỏng vấn 30 người về một bộ phim mới chiếu trên truyền hình. Người điều tra yêu cầu cho điểm bộ phim (thang điểm là 100). Kết quả được trình bày trong bảng phân bố tần số ghép lớp sau đây
	\begin{center}
	\begin{tabular}{|c|c|c|c|c|c|}
	\hline
	Số điểm  & $[50;60)$ & $[60; 70)$ & $[70;80)$ & $[80;90)$ & $[90;100)$ \\ \hline
	Số người & 2           & 6          & 10          & 8           & 4           \\ \hline
	\end{tabular}
	\end{center}
	\choiceTF
	{\True Ước lượng số trung bình của mẫu ghép là $77$} 
	{\True Giá trị đại diện của nhóm ${[90 ; 100)}$ là $95$}
	{Nhóm chứa mốt của mẫu số liệu trên là nhóm $[80;90)$}
	{Mốt của mẫu số liệu là $74{,}67$}
	\loigiai{
		\begin{itemchoice}
		\itemch \textbf{Đúng}. 
				\begin{center}
				\begin{tabular}{|c|c|c|c|c|c|}
				\hline
				Số điểm  & $[50;60)$ & $[60; 70)$ & $[70;80)$ & $[80;90)$ & $[90;100)$ \\ \hline
				Giá trị đại diện & $55$           & $65$          & $75$         & $85$           & $95$          \\ \hline
				Số người & 2           & 6          & 10          & 8           & 4           \\ \hline
				\end{tabular}
				\end{center}
		Ước lượng số trung bình của mẫu ghép là:
		$$\bar{x}=\dfrac{55\cdot2+65\cdot6+75\cdot10+85\cdot8+95\cdot4}{30}=77.$$
		\itemch \textbf{Đúng}.Giá trị đại diện của nhóm ${[90 ; 100)}$ là $95$.
		\itemch \textbf{Sai}. Nhóm chứa mốt của mẫu số liệu trên là nhóm $[70; 80)$.
		\itemch \textbf{Sai}. 
		Do đó, $u_m=70; n_{m-1}=6 ; n_m=10 ; n_{m+1}=8 ; u_{m+1}-u_m=80-70=10$.\\
		$M_{\circ}=70+\dfrac{10-6}{(10-6)+(10-8)} \cdot 10=76{,}67.$
		\end{itemchoice}
		}
\end{ex}

\Closesolutionfile{ans}
\indapan{2}{ans/ans-11-C5B1-DE2-DS}

\caukq
\Opensolutionfile{ans}[ans/ans-11-C5B1-DE2-KQ]
\begin{ex}%[1D5H1-3]
	Mẫu số liệu sau ghi lại cân nặng của $30$ bạn học sinh (đơn vị: kilôgam)
	\[
		\begin{array}{cccccccccc}
			17 & 40 & 39 & 40{,}5 & 42 & 51 & 41{,}5 & 39 & 41 & 30\\
			40 & 42 & 40{,}5 & 39{,}5 & 41 & 40{,}5 & 37 & 39{,}5 & 40 & 41\\
			38{,}5 & 39{,}5 & 40 & 41 & 39 & 40{,}5 & 40 & 38{,}5 & 39{,}5 & 41{,}5
		\end{array}
	\]
	Xác định số trung bình cộng của mẫu số liệu ghép nhóm trên.
	\shortans{$40$}
	\loigiai{
				Trung bình cộng của mẫu số liệu trên là
				\[
				\overline{x} = \dfrac{17{,}5 \cdot 1 + 32{,}5 \cdot 1 + 37{,}5 \cdot 10+ 42{,}5 \cdot 17+ 52{,}5 \cdot 1}{30} = 	40\text{ (kg)}.
				\]
	}
\end{ex}

\begin{ex}%[1D5H1-3]
	\immini{
		Bảng bên cho ta bảng tần số ghép nhóm số liệu thống kê chiều cao của $40$ mẫu cây ở một vườn thực vật (đơn vị: centimét).
		Xác định số trung bình cộng của mẫu số liệu ghép nhóm trên.
	}
	{
		\begin{tabular}{|c|c|c|}
			\hline
			\textbf{Nhóm} & \textbf{Tần số} & \textbf{Tần số tích luỹ}\\ 
			\hline
			$\left[30;40\right)$ & $4$ & $4$\\
			$\left[40;50\right)$ & $10$ & $14$\\
			$\left[50;60\right)$ & $14$ & $28$\\
			$\left[60;70\right)$ & $6$ & $34$\\
			$\left[70;80\right)$ & $4$ & $38$\\
			$\left[80;90\right)$ & $2$ & $40$\\
			\hline
			& $n = 40$ &\\
			\hline
		\end{tabular}
	}
	\shortans{$55{,}5$}
\loigiai{
		Trung bình cộng của mẫu số liệu trên là
			\[
			\overline{x} = \dfrac{35 \cdot 4 + 45 \cdot 10 + 55 \cdot 14 + 65 \cdot 6 + 75 \cdot 4 + 85 \cdot 2}{40} = 	55{,}5\text{ (cm)}.
			\]
}
\end{ex}

\begin{ex}%[1D5H1-4]
	Kết quả kiểm tra môn Toán của lớp $11D$ như sau
	\[
		\begin{array}{cccccccccccccccccccc}
			5 & 6 & 7 & 5 & 6 & 9 & 10 & 8 & 5 & 5 & 4 & 5 & 4 & 5 & 7 & 4 & 5 & 8 & 9 & 10 \\
			5 & 3 & 5 & 6 & 5 & 7 & 5 & 8 & 4 & 9 & 5 & 6 & 5 & 6 & 8 & 8 & 7 & 9 & 7 & 9
		\end{array}
	\]
	Mốt của bảng số liệu ghép nhóm trên là bao nhiêu (làm tròn kết quả đến hàng phần mười)?
	\shortans{$6{,}2$}
\loigiai{\immini{Ta thấy: Nhóm $2$ ứng với nửa khoảng $\left[5;7\right)$ là nhóm có tần số lớn nhất với $u=5$, $g=2$, $n_2 = 18$. Nhóm $1$ có tần số $n_1 = 5$, nhóm $3$ có tần số $n_3=10$.\\
			Khi đó, mốt của mẫu số liệu là 
			\[
				M_o = 5 + \left( \dfrac{18- 5}{2\cdot 18 - 5 - 10} \right) \cdot 2 \approx 6{,}2.
			\]
	}{
		\begin{tabular}{|c|c|}
			\hline
			\textbf{Nhóm} & \textbf{Tần số}\\ 
			\hline
			$\left[3;5\right)$ & $5$\\
			\hline
			$\left[5;7\right)$ & $18$\\
			\hline
			$\left[7;9\right)$ & $10$\\
			\hline
			$\left[9;11\right)$ & $7$\\
			\hline
			& $n = 40$ \\
			\hline
		\end{tabular}
	}
}
\end{ex}

\begin{ex}%[1D5H1-3]
	\immini{Một nhà thực vật học đo chiều dài của $74$ lá cây (đơn vị: milimét) và thu được bảng tần số như bảng bên. Tính chiều dài trung bình của $74$ lá cây trên theo đơn vị milimét (làm tròn kết quả đến hàng phần trăm).}
	{
		\begin{tabular}{|c|c|c|}
			\hline
			\textbf{Nhóm} & \textbf{Giá trị đại diện} & \textbf{Tần số}\\ 
			\hline
			$\left[5{,}45;5{,}85\right)$ & $5{,}65$ & $5$\\
			$\left[5{,}85;6{,}25\right)$ & $6{,}05$ & $9$\\
			$\left[6{,}25;6{,}65\right)$ & $6{,}45$ & $15$\\
			$\left[6{,}65;7{,}05\right)$ & $6{,}85$ & $19$\\
			$\left[7{,}05;7{,}45\right)$ & $7{,}25$ & $16$\\
			$\left[7{,}45;7{,}85\right)$ & $7{,}65$ & $8$\\
			$\left[7{,}85;8{,}25\right)$ & $8{,}05$ & $2$\\
			\hline
			&  & $n = 74$\\
			\hline
		\end{tabular}
	}
	\shortans{$6{,}80$}
	\loigiai{
		Chiều dài trung bình của $74$ lá cây mà nhà thực vật học đo xấp xỉ là 
		\[
			\overline{x} = \dfrac{5\cdot 5{,}65 + 9 \cdot 6{,}05 + 15\cdot 6{,}45 + 19\cdot 6{,}85 + 16 \cdot 7{,}25 + 8\cdot 7{,}65 + 2\cdot 8{,}05}{74} \approx 6{,}80\ (\text{mm}).
		\]
	}
\end{ex}

\begin{ex}%[1D5H1-4]
		Người ta ghi lại tuổi thọ của một số con ong cho kết quả như sau
		\begin{center}
			\begin{tabular}{|c|c|c|c|c|c|}
				\hline 
				Tuổi thọ (ngày) & $[0;20)$ &$[20;40)$&$[40;60)$&$[60;80)$&  $[80;100)$\\
				\hline 
				Số lượng&$5$&$12$&$23$&$31$&$29$\\
				\hline 
			\end{tabular}
		\end{center}
	Tìm mốt của mẫu số liệu.
	\shortans{$76$}
	\loigiai{
	Tần số lớn nhất là $31$ nên nhóm chứa mốt là nhóm $[60;80)$.\\
	Ta có $j=4$; $a_4=60$; $m_4=31$; $m_3=23$; $m_5=29$; $h=20$.\\
	Khi đó $M_O=60+\dfrac{31-23}{(31-23)+(31-29)}\cdot 20=76$.
	} 
\end{ex}

\begin{ex}%[1D5H1-3]
	Cơ cấu dân số Việt Nam năm 2020 theo độ tuổi được cho trong bảng sau 
	\begin{center}
		\begin{tabular}{|c|c|c|c|c|c|}
			\hline 
			Độ tuổi & Dưới $5$ tuổi &$5-14$&$15-24$&$25-64$& Trên $65$\\
			\hline 
			Số người (triệu)&$7{,}89$&$14{,}68$&$13{,}32$&$53{,}78$&$7{,}66$\\
			\hline 
		\end{tabular}
	\end{center} 		
	Chọn $80$ là giá trị đại diện cho nhóm trên $65$ tuổi. Tính tuổi trung bình của người Việt Nam năm $2020$. Làm tròn đến hàng phần mười.
	\shortans{$35{,}2$}
	\loigiai{
	Trong mỗi khoảng độ tuổi, giá trị đại diện là trung bình cộng của giá trị hai đầu mút nên ta có bảng sau
	\begin{center}
		\begin{tabular}{|c|c|c|c|c|c|}
			\hline 
			Độ tuổi & $2{,}5$ &$9{,}5$&$19{,}5$&$44{,}5$& $80$\\
			\hline 
			Số người (triệu)&$7{,}89$&$14{,}68$&$13{,}32$&$53{,}78$&$7{,}66$\\
			\hline 
		\end{tabular}
	\end{center} 
Tổng số dân là $n=97{,}33$ triệu người. Do đó tuổi trung bình của người dân Việt Nam trong năm $2020$ là
$$\overline{x}=\dfrac{2{,}5\cdot 7{,}89+9{,}5\cdot 14{,}68+19{,}5\cdot 13{,}32+44{,}5\cdot 53{,}78+80\cdot 7{,}66}{97{,}33}\approx 35{,}2.$$
	} 
\end{ex}
\Closesolutionfile{ans}
\indapan{6}{ans/ans-11-C5B1-DE2-KQ}
